% Chapter 17, generated by make_chapters.py from papers/Ch17_Bio_11_Isoform_Polytope/anccft24.tex.
% Source paper: Bio 11
\chapter[The isoform decomposition polytope]{The isoform decomposition polytope of a splicing graph, and a BEST identity for its integer points}\label{chap:ch17}
\begin{chaptersummary}
Junction and exon read counts of a gene are an integer flow $w$ on its splicing graph, a
directed acyclic graph from the transcription start to the polyadenylation site, and the
transcript isoforms are $s$--$t$ paths. The abundances compatible with $w$ form the fibre
polytope $\Th(w)=\{\theta\ge0:\sum_j\theta_j\mathbf 1_{P_j}=w\}$. We prove that its dimension
is the number of $s$--$t$ paths in the support of $w$ minus $|E_w|-|V_w|+2$, because $w$
lies in the relative interior of the cone spanned by those paths and the path--arc matrix has
the rank of the flow space; for alternative events in series with $k_1,\dots,k_r$ branches
this is $\prod k_i-\sum(k_i-1)-1$, so deconvolution from junction counts is unique for one
event and fails already for two cassette exons, and for Dscam ($12\times48\times33\times2$)
the $38\,016$ isoforms sit in a $37\,924$-dimensional polytope. Every minimum-support
decomposition, with real or integer weights, is a vertex of $\Th(w)$; but the vertices of
$\Th(w)$ are not integral in general --- three parallel-arc pairs in series and $w\equiv1$
have a vertex with four weights $\tfrac12$ beside an integer minimum of two paths --- so the
extreme rays of the flow cone, the vertices of the fibre, and integrality are three different
things. The new algebraic content is an identity. Adding $F$ return arcs $t\to s$ to the
multigraph $G_w$, $F$ the flow value, the per-sequence BEST count of the augmented graph
equals $\sum_\theta (F-1)!/\prod_j\theta_j!$ over the integer points of $\Th(w)$: a
determinant computes the multinomial-weighted partition function of the integer
decompositions, of which the plain number of integer points has no such expression. The
identity and the vertex statements are verified exactly on $238$ flows, and no instance with
a real minimum below the integer minimum was found among $232$ random small flows.
\end{chaptersummary}

\medskip
\noindent\textbf{Keywords.} splicing graph, alternative splicing, flow decomposition, isoform
quantification, identifiability, BEST theorem, fibre polytope, Dscam.

\noindent\textbf{MSC 2020.} 05C20, 05C21, 05C45, 52B12, 90C27, 92D20.

\section{What was asked, and what is true}

The Bio strand of these notes reads assembly ambiguity off the critical group of a de Bruijn
graph Chapter~\ref{chap:ch01}, double-stranded ambiguity off the box points of an eigenlattice of the
strand involution Chapter~\ref{chap:ch07}, and polyploid phasing off the $p$-fold decompositions of a
pooled flow Chapter~\ref{chap:ch13}. The common object is the decomposition of a balanced integer flow
into closed walks. Transcript quantification is the same object with the circle cut open: a
splicing graph is a DAG with a source and a sink \cite{HAST02}, read counts are a flow, and
the isoforms are $s$--$t$ paths; the question ``which isoforms, in what abundance'' is the
question of decomposing the flow into paths, and the ambiguity is the polytope of such
decompositions.

The specification of this note asked for a proof that ``minimum flow decompositions
correspond to convex combinations of integer lattice points on the extreme rays of the flow
polytope''. That sentence runs three distinct objects together, and it is worth separating
them because two of the three are elementary and the third is false as stated. The extreme
rays of the cone of $s$--$t$ flows of a DAG are the path indicators, and the vertices of the
polytope of flows of fixed value are the paths scaled by the value: this is the total
unimodularity of the incidence matrix and nothing else. The polytope that matters for a
\emph{given} $w$ is the fibre $\Th(w)$ of the path--arc map over $w$, whose vertices are the
decompositions with linearly independent support; every minimum-support decomposition is one
of them (Proposition~\ref{ch17:prop:vertex}). But the path--arc matrix is not totally unimodular,
and $\Th(w)$ has fractional vertices for integer $w$ (Proposition~\ref{ch17:prop:frac}); the
integer decompositions are the lattice points of $\Th(w)$, of which the vertices are in
general not a subset. Minimum flow decomposition is NP-hard \cite{HHKS12}, and nothing here
changes that.

What is true and, we believe, new is in Sections~\ref{ch17:sec:dim} and~\ref{ch17:sec:best}. The
dimension of $\Th(w)$ is exactly the number of paths in the support of $w$ minus the rank of
the flow space of that support (Theorem~\ref{ch17:thm:dim}); this is the precise form of the
non-identifiability of isoform deconvolution from junction counts, a phenomenon known since
\cite{LSGB08,HJXW09}, and it gives the closed form $\prod k_i-\sum(k_i-1)-1$ for events in
series and the number $37\,924$ for Dscam. And the BEST theorem, applied to the flow closed
up by $F$ return arcs, evaluates the multinomial-weighted sum over the integer points of
$\Th(w)$ as a determinant (Theorem~\ref{ch17:thm:best}): the same normalisation by
$\prod c(a)!$ that Chapter~\ref{chap:ch07} needed for double strands appears here as the passage from
cyclic sequences of isoforms to abundance vectors. The script
\texttt{isoform\_polytope.py} verifies every statement and reproduces every number; it prints
a seven-line pass/fail battery.

\section{The splicing graph, its flows, and the fibre polytope}\label{ch17:sec:dim}

\begin{definition}
A \emph{splicing graph} is a directed acyclic multigraph $G=(V,E)$ with a source $s$ and a
sink $t$ in which every arc lies on some $s$--$t$ path. An \emph{isoform} is an $s$--$t$
path; $\mathcal P$ is the set of isoforms and $A\in\{0,1\}^{E\times\mathcal P}$ the
path--arc matrix. A \emph{flow} is $w\in\Z^E_{\ge0}$ with $\partial w=F(\delta_t-\delta_s)$,
$F$ its \emph{value}; $E_w=\supp w$, $V_w$ the vertices touched, $\mathcal P(w)$ the isoforms
inside $E_w$, and
\[
  \Th(w)=\Big\{\theta\in\R_{\ge0}^{\mathcal P(w)}:\ \sum_{P\in\mathcal P(w)}\theta_P\mathbf 1_P=w\Big\}.
\]
\end{definition}

$\Th(w)$ is non-empty by flow decomposition (a non-negative flow on a DAG is a non-negative
combination of $s$--$t$ paths, and an integer flow of integer paths), and bounded since
$\theta_P\le F$. It is the polytope of abundance vectors compatible with the counts.

\begin{lemma}\label{ch17:lem:rank}
$\rank A=|E|-|V|+2$.
\end{lemma}

\begin{proof}
The columns of $A$ lie in $\mathcal F=\{x\in\Q^E:\partial x\in\Q(\delta_t-\delta_s)\}$,
which has dimension $\dim\ker\partial+1=|E|-|V|+2$ since $G$ is connected. They span it:
let $x\in\mathcal F$ and let $x_0=\sum_P\mathbf 1_P$, which is positive on every arc because
every arc lies on a path; for $M$ large, $x+Mx_0$ is a non-negative flow, hence a
non-negative combination of paths, so $x$ is in the span.
\end{proof}

\begin{theorem}[Dimension of the isoform polytope]\label{ch17:thm:dim}
$\dim\Th(w)=|\mathcal P(w)|-\big(|E_w|-|V_w|+2\big)$. In particular the abundances are
determined by the counts if and only if the isoforms inside the support of $w$ are linearly
independent as arc vectors, i.e.\ $|\mathcal P(w)|=|E_w|-|V_w|+2$.
\end{theorem}

\begin{proof}
The subgraph $G_w=(V_w,E_w)$ is a splicing graph in its own right (every arc of positive flow
lies on a path of positive flow, by flow decomposition), so Lemma~\ref{ch17:lem:rank} gives
$\rank A_w=|E_w|-|V_w|+2$ for its path--arc matrix. The cone $C_w=\{x\in\mathcal F_w:x\ge0\}$
is the cone generated by $\mathcal P(w)$, and its relative interior is $\{x\in\mathcal
F_w:x_a>0\ \forall a\in E_w\}$, because every coordinate $x_a$ is a non-trivial linear form
on $\mathcal F_w$ (some path uses $a$). Since $w>0$ on $E_w$, $w\in\relint C_w$, and a point
in the relative interior of a finitely generated cone is a strictly positive combination of
all the generators. So $\Th(w)$ contains a point with full support, its affine hull is the
whole affine space $\{\theta:A_w\theta=w\}$, and the dimension is $|\mathcal P(w)|-\rank A_w$.
\end{proof}

\begin{corollary}[Events in series]\label{ch17:cor:series}
If $G$ is a chain of $r$ alternative events, the $i$-th offering $k_i$ branches between
consecutive junctions, then for a flow of full support
\[
  \dim\Th(w)=\prod_{i=1}^r k_i-\sum_{i=1}^r(k_i-1)-1 ,
\]
which is $0$ if and only if at most one $k_i$ exceeds $1$. Two cassette exons in series give
$4-2-1=1$; three give $4$; a cassette, two pairs of mutually exclusive exons, a second
cassette and a four-way mutually exclusive block give $64-8=56$; the four blocks of
\emph{Dscam} \cite{SCH00}, with $12,48,33,2$ alternatives, give $38\,016-92=37\,924$.
\end{corollary}

\begin{proof}
Paths multiply, $|\mathcal P|=\prod k_i$; branches with $\ell$ arcs have $\ell-1$ interior
vertices, so $|E|-|V|+2=\sum_i k_i-r-1+2=\sum_i(k_i-1)+1$ independently of the branch
lengths. The last claim is Theorem~\ref{ch17:thm:dim}.
\end{proof}

The identifiability statement in Theorem~\ref{ch17:thm:dim} is the rank condition of
\cite{HJXW09} specialised to junction-and-exon counts; the dimension formula makes it
quantitative, and Corollary~\ref{ch17:cor:series} says that for any gene with two independent
alternative events the counts do not determine the abundances, whatever the depth of
sequencing. Table~\ref{ch17:tab:models} lists the models and their dimensions.

\section{Vertices, minimum decompositions, and integrality}\label{ch17:sec:vertices}

\begin{proposition}\label{ch17:prop:vertex}
The vertices of $\Th(w)$ are the decompositions whose support is a linearly independent set
of path vectors. Every decomposition of minimum support --- among real decompositions, or
among integer ones --- is a vertex.
\end{proposition}

\begin{proof}
The first sentence is the characterisation of basic feasible solutions of
$\{\theta\ge0:A\theta=w\}$. For the second, let $\theta$ have minimum support and suppose the
columns of its support are dependent, $Av=0$ with $v\ne0$ supported in $\supp\theta$. Then
$\theta+\tau v\in\Th(w)$ for $\tau$ in a closed interval containing $0$ in its interior, and
at an endpoint of that interval some coordinate of $\supp\theta$ vanishes, giving a
decomposition of smaller support; if $\theta$ and $v$ are integral, the endpoints can be
taken at integer $\tau$ (the first coordinate to vanish does so at $\tau=-\theta_j/v_j$, and
one may stop at the last integer before it, or note that $v$ may be scaled to make the
endpoint integral without changing the argument). Contradiction.
\end{proof}

\begin{proposition}[A fractional vertex]\label{ch17:prop:frac}
Let $G$ consist of four vertices $s=v_0,v_1,v_2,v_3=t$ with two parallel arcs
$a_i,a_i'$ from $v_{i-1}$ to $v_i$, and $w\equiv1$ (value $2$). Then $\Th(w)$ has $8$
isoforms, dimension $4$, four integer points, and six vertices of which two are fractional:
$\theta=\tfrac12$ on the four paths $(a_1,a_2',a_3)$, $(a_1,a_2,a_3')$, $(a_1',a_2,a_3)$,
$(a_1',a_2',a_3')$ is a vertex, while the minimum decompositions are the two integer
points with two paths each.
\end{proposition}

\begin{proof}
The four listed paths cover every arc exactly twice, so $\theta$ is feasible; their arc
vectors are independent, because a relation $\alpha P_1+\beta P_2+\gamma P_3+\delta P_4=0$
read on $a_1,a_2,a_3$ gives $\alpha+\beta=\beta+\gamma=\alpha+\gamma=0$, hence
$\alpha=\beta=\gamma=0$ and then $\delta=0$. So $\theta$ is a vertex by
Proposition~\ref{ch17:prop:vertex}, and it is not integral. The $3\times3$ submatrix of the
path--arc matrix on those three arcs and three paths has determinant $2$, which is where the
halves come from. The counts $8,4,4,6,2$ are the transcript of the script.
\end{proof}

So the lattice points of $\Th(w)$ and its vertices are different sets, each containing the
minimum decompositions. Whether the integer minimum can exceed the real minimum we do not
know: the script compared the two on $232$ random flows on small DAGs with $3$ to $12$
isoforms and found no instance where they differ, and one instance with a fractional
vertex. Minimum flow decomposition being NP-hard \cite{HHKS12}, with the bound
$|E|-|V|+2$ on the number of paths needed \cite{VCCM08} and exact solvers in practice
\cite{KKO18,DT22}, nothing in this section is intended as an algorithm.

\section{The BEST identity}\label{ch17:sec:best}

Close the flow up. Let $\Ghat_w$ be the multigraph with $w(a)$ labelled parallel copies of
each arc $a\in E_w$ and $F$ labelled \emph{return arcs} $t\to s$. It is balanced and
connected, so it has Eulerian circuits, and the per-sequence count of Chapter~\ref{chap:ch07} applies:
\[
  \Nseq(\Ghat_w)=\frac{t(\Ghat_w)\,\Loc(\Ghat_w)}{F!\,\prod_{a}w(a)!}
  =\sum_{T}\frac1{d(T)},
\]
the sum over the circular sequences $T$ of $\Ghat_w$ --- cyclic sequences of arc contents
using each content with its multiplicity --- weighted by the order $d(T)$ of the rotation
group of $T$; here $t$ is the number of spanning arborescences converging on a root and
$\Loc=\prod_v(d^+(v)-1)!$.

\begin{theorem}[BEST identity for isoform decompositions]\label{ch17:thm:best}
\[
  \Nseq(\Ghat_w)\;=\;\sum_{\theta\in\Th(w)\cap\Z^{\mathcal P(w)}}\ \frac{(F-1)!}{\prod_{P}\theta_P!}\,.
\]
\end{theorem}

\begin{proof}
A circular sequence of $\Ghat_w$ uses the $F$ return arcs, and between two consecutive
returns it walks from $s$ to $t$ along arcs of the DAG; a walk in a DAG repeats no vertex,
so it is an isoform. Thus $T$ is a cyclic sequence $(P_{j_1},\dots,P_{j_F})$ of isoforms,
its arc usage is $\sum_i\mathbf 1_{P_{j_i}}=w$, and the multiplicity vector
$\theta_P=\#\{i:P_{j_i}=P\}$ is an integer point of $\Th(w)$ with $\sum_P\theta_P=F$.
Conversely every integer point $\theta$ and every cyclic arrangement of the multiset
$\{P^{\theta_P}\}$ gives such a $T$. For a fixed $\theta$ the cyclic arrangements, weighted by
$1/d(T)$, number $\frac1F\cdot\frac{F!}{\prod_P\theta_P!}$: the cyclic group of order $F$
acts on the $F!/\prod\theta_P!$ linear arrangements, and $\sum_{\text{orbits}}1/|\mathrm{stab}|
=|\text{set}|/|\text{group}|$.
\end{proof}

The left side is a determinant times factorials; the right side is a sum over the lattice
points of a polytope of the multinomial coefficient of each point, i.e.\ the number of
linear orders of $F$ transcripts with those abundances. Theorem~\ref{ch17:thm:best} is thus a
closed form for the ``partition function'' of the integer decompositions with the natural
combinatorial weight, and it is the isoform analogue of the fact that BEST counts genome
reconstructions: there the whole spectrum is one closed walk, here it is a cyclic word in
the isoforms. The plain number of integer points $|\Th(w)\cap\Z^{\mathcal P}|$, which is
what a uniform prior over decompositions would need, is a lattice-point count with no
determinantal formula, and Theorem~\ref{ch17:thm:best} says precisely which weighting is
determinantal.

\begin{remark}
On one strand, $\prod w(a)!$ is a constant and the single-genome notes count arc-rooted
circuits without it. Here, as in the double-stranded and polyploid sums, the divisor varies
with the point: for the tropomyosin-like flow of Table~\ref{ch17:tab:flows} the arborescence count
is $t=3\,628\,800\,000$, and only after multiplying by $\Loc$ and dividing by
$F!\prod_aw(a)!$ does one obtain $\Nseq=22\,680\,000$, the quantity that matches the lattice
sum.
\end{remark}

\section{Verification and data}\label{ch17:sec:data}

\begin{table}[t]
\centering
\small
\begin{tabular}{lrrrrrl}
\toprule
model & $|V|$ & $|E|$ & isoforms & $|E|-|V|+2$ & $\dim\Th$ & unique?\\
\midrule
cassette exon (skipping)                 &   3 &   3 &      2 &  2 &     0 & yes\\
mutually exclusive exon pair             &   4 &   4 &      2 &  2 &     0 & yes\\
alternative $5'$ splice site             &   2 &   2 &      2 &  2 &     0 & yes\\
two cassettes in series                  &   5 &   6 &      4 &  3 &     1 & no\\
three cassettes in series                &   7 &   9 &      8 &  4 &     4 & no\\
tropomyosin-like (five events)           &  16 &  22 &     64 &  8 &    56 & no\\
three parallel-arc pairs                 &   4 &   6 &      8 &  4 &     4 & no\\
\emph{Dscam} ($12\times48\times33\times2$) & 100 & 190 & 38\,016 & 92 & 37\,924 & no\\
\bottomrule
\end{tabular}
\caption{Ranks and dimensions. The tropomyosin-like model is a pair of mutually exclusive
exons, a cassette, a second pair, a second cassette and a four-way mutually exclusive block
in series. The rank of the path--arc matrix equals $|E|-|V|+2$ in every row (over $\Q$ for
the small models, modulo $2^{31}-1$ for the $190\times38\,016$ Dscam matrix, $12$\,s).
Checks (R), (D), (U), (S).}
\label{ch17:tab:models}
\end{table}

\begin{table}[t]
\centering
\small
\begin{tabular}{lrrrrrrrrr}
\toprule
flow & $F$ & paths & $\dim$ & int.\ pts & vert. & frac. & MFD & $t$ & $\Nseq$\\
\midrule
two cassettes      &  6 &  4 &  1 &   3 &   2 &  0 & 3 & 216 & 50\\
three parallel pairs &  2 &  8 &  4 &   4 &   6 &  2 & 2 & 8 & 4\\
tropomyosin-like   & 10 & 24 & 18 & 714 & 358 & 90 & 4 & $3.6288\cdot10^{9}$ & $22\,680\,000$\\
three cassettes    &  7 &  4 &  1 &   2 &   2 &  0 & 3 & 57\,624 & 35\\
\bottomrule
\end{tabular}
\caption{Flows: two cassettes with $w=(3,3,3,2,4,2)$; three parallel pairs with $w\equiv1$;
the tropomyosin-like model with a flow generated from five random isoforms of weights $1$ to
$3$; three cassettes with a random flow. ``paths'' is $|\mathcal P(w)|$; ``MFD'' the minimum
support, equal for real and integer weights in every row; $t$ the arborescence count of
$\Ghat_w$; $\Nseq$ its per-sequence count, which equals $\sum_\theta(F-1)!/\prod\theta_P!$
over the integer points in every row. In the tropomyosin-like row the flow was generated from
five isoforms and the minimum decomposition uses four. Checks (V), (B).}
\label{ch17:tab:flows}
\end{table}

\paragraph{The seven checks.} (R) $\rank A=|E|-|V|+2$ on the seven small models; (D) the
dimension formula and its series form; (U) uniqueness for one event, dimension one for two
cassettes; (S) Dscam; (V) every minimum-support integer decomposition is a vertex, and the
fractional vertex of Proposition~\ref{ch17:prop:frac} exists; (B) the BEST identity on the four
flows of Table~\ref{ch17:tab:flows} and on $232$ random flows on DAGs with $4$ to $6$ vertices;
(I) real versus integer minimum on those $232$ flows --- never different --- with one
fractional vertex among them. All pass; the run takes about three and a half minutes, almost
all of it in the random search.

\section{Scope, provenance, and what is open}

\paragraph{Scope.} Lemma~\ref{ch17:lem:rank}, Theorem~\ref{ch17:thm:dim},
Corollary~\ref{ch17:cor:series}, Proposition~\ref{ch17:prop:vertex} and Theorem~\ref{ch17:thm:best} are
proved for arbitrary splicing graphs; Proposition~\ref{ch17:prop:frac} is one graph. The model is
the exact flow: no read-length effects, no paired reads, no noise, and no length
normalisation of abundances; reads spanning several junctions are additional linear
constraints that cut $\Th(w)$ by hyperplanes and lower its dimension, which is how phasing
information enters \cite{HJXW09}. Minimum flow decomposition is treated only to the extent
of Section~\ref{ch17:sec:vertices}.

\paragraph{Provenance.} Splicing graphs are \cite{HAST02}; the non-identifiability of
isoform deconvolution is \cite{LSGB08,HJXW09}; the $|E|-|V|+2$ bound on the number of paths
in a decomposition and the greedy algorithms are \cite{VCCM08}; NP-hardness of minimum flow
decomposition is \cite{HHKS12}; practical exact solvers are \cite{KKO18,DT22}; the Dscam
isoform count is \cite{SCH00}; the BEST theorem is \cite{BEST,SmithTutte} and its
per-sequence normalisation is Chapter~\ref{chap:ch07}. The tropical geometry of statistical models,
which the planning of this note once proposed to bring in, is \cite{PS04} and is not used:
nothing in the polytope $\Th(w)$ is tropical. We are not aware of prior work stating the
dimension of the isoform fibre in closed form, or of the BEST identity of
Theorem~\ref{ch17:thm:best}; a targeted search found neither.

\paragraph{Open.} \emph{First}, whether the integer minimum can exceed the real minimum
for an integer flow on a DAG. \emph{Second}, a formula or complexity classification for the
number of integer points of $\Th(w)$, as opposed to the multinomial-weighted sum of
Theorem~\ref{ch17:thm:best}. \emph{Third}, the effect of multi-junction reads on the dimension:
each read spanning a set $J$ of junctions with allele pattern $\pi$ adds the constraint that
the total abundance of isoforms with pattern $\pi$ on $J$ be the observed count, and the
dimension after all such constraints is the true residual ambiguity of a sequencing design.
\emph{Fourth}, the analogue of Theorem~\ref{ch17:thm:best} with the return arcs weighted, which
would give a generating function in $F$ rather than a single number.
